This section lists several functions, definitions, etc., handy for the GCA lab. It is a living document, where updates are added according to the lab sessions and feedback from the students. Thanks for pointing things out.

\textbf{GPU memory allocation}
\begin{lstlisting}[language=CUDA,numbers=none]
float* d_A;
cudaMalloc(&d_A, size);
\end{lstlisting}

\textbf{Memory migration to / from GPU}
\begin{lstlisting}[language=CUDA,numbers=none]
float* h_A = (float*)malloc(size);
cudaMemcpy(d_A, h_A, size, cudaMemcpyHostToDevice);
cudaMemcpy(h_A, d_A, size, cudaMemcpyDeviceToHost);
\end{lstlisting}

\textbf{Kernel invocation (100 blocks, with 80 threads each)}
\begin{lstlisting}[language=CUDA,numbers=none]
foo<<<100, 80>>>(d_A);
\end{lstlisting}

\textbf{Thread indexing}
\begin{lstlisting}[language=CUDA,numbers=none]
int i = threadIdx.x;
\end{lstlisting}

\textbf{Setting to global scope (invoked by CPU or GPU)}
\begin{lstlisting}[language=CUDA,numbers=none]
__global__ void foo(d_A){...}
\end{lstlisting}

\textbf{Setting to GPU scope (invoked by GPU)}
\begin{lstlisting}[language=CUDA,numbers=none]
__device__ float foo(d_A){...}
\end{lstlisting}

\textbf{Timing (CPU)}
\begin{lstlisting}[language=CUDA,numbers=none]
#include <chrono>
const auto start = std::chrono::steady_clock::now();
...
const auto end = std::chrono::steady_clock::now();
const std::chrono::duration<double> elapsed_seconds{end - start};
\end{lstlisting}

\textbf{Timing (GPU)}
\begin{lstlisting}[language=CUDA,numbers=none]
cudaEvent_t start, stop;
cudaEventCreate(&start);
cudaEventCreate(&stop);
cudaEventRecord(start);
...
cudaEventSynchronize();
cudaEventRecord(stop);
float milliseconds;
cudaEventElapsedTime(&milliseconds, start_cuda, stop_cuda);
\end{lstlisting}
